%% chapters/01_introduction.tex

\section{Introduction}
\label{sec:intro}

Fraudulent transactions impose enormous costs on the global financial system.
Industry estimates suggest that card fraud alone exceeds \ billion annually,
with e-commerce fraud growing at compound rates above 20\% \citep{ieeecisFraud2019}.
Traditional rule-based detection systems suffer from two fundamental limitations:
they require expert-crafted rules that become stale as fraud patterns evolve,
and they cannot express the probabilistic uncertainty inherent to anomaly scoring.

Machine learning approaches address these limitations but introduce their own
challenges. Fraud data is heavily imbalanced --- legitimate transactions
constitute 96--99\% of all records --- which invalidates accuracy as a metric
and demands specialised resampling or cost-sensitive training strategies.
Furthermore, a deployed model must not merely classify; it must provide
a defensible, quantified cost-benefit trade-off that allows an analyst or
compliance team to set the operating threshold rationally.

This report describes \textbf{ARGUS} (Adaptive Real-time Grading and
Unsupervised Scoring), a complete seven-layer fraud detection system
structured to mirror the architecture of a real data science or trading
analytics team:

\begin{enumerate}
    \item \textbf{Ingestion:} IEEE-CIS transaction and identity data loaded
          into a normalised SQLite schema.
    \item \textbf{Storage:} Relational tables with referential integrity
          and indexed query paths.
    \item \textbf{Processing:} A leak-free \texttt{sklearn} Pipeline covering
          log-amount features, cyclical time encoding, missing-value flags,
          frequency encoding, leak-free target encoding, and PCA compression
          of 339 Vesta V-features.
    \item \textbf{Modelling:} Four models spanning supervised and unsupervised
          paradigms, with Bayesian hyperparameter optimisation and rigorous
          cross-validation.
    \item \textbf{Evaluation:} AUPRC-primary benchmarking, bootstrap
          confidence intervals, Wilcoxon significance tests, and cost-curve
          analysis for optimal threshold selection.
    \item \textbf{Serving:} FastAPI endpoint containerised with Docker,
          with real-time SHAP explanations and {99} < 100$\,ms latency.
    \item \textbf{Automation:} Apache Airflow DAGs for weekly ingestion and
          AUPRC-gated model retraining.
\end{enumerate}

\subsection{Contributions}

The primary contributions of this work are:
\begin{itemize}
    \item A rigorous benchmark of four fraud detection models (supervised and
          unsupervised) on the IEEE-CIS dataset with statistical confidence intervals.
    \item A cost-theoretic framework for decision threshold selection that
          replaces the arbitrary $\theta = 0.5$ convention with Expected Cost
          minimisation.
    \item A fully deployable, containerised inference pipeline with
          Prometheus-format monitoring, SHAP explainability, and Population
          Stability Index drift detection.
    \item Complete reproducibility: all code, experiments, and this report
          are version-controlled and public at
          \url{https://github.com/rudraakshreddy/argus}.
\end{itemize}
